\documentclass[10pt,a4paper]{amsart}
\usepackage[margin=19mm]{geometry}
\usepackage{amsmath,amssymb,mathtools}
\usepackage[scr=rsfs,cal=euler]{mathalfa}
\usepackage{xfrac}
\usepackage[unicode,hidelinks,bookmarks=false]{hyperref}
\newcommand{\ie}{i.e.\ }
\newcommand{\cf}{cf.\ }
\newcommand{\resp}{respectively\ }
\newcommand{\Subsection}{\subsection}
\providecommand{\nobreakdash}{\nobreak\hbox{-}}
\setlength{\emergencystretch}{2em}
\multlinegap0pt
\usepackage{amssymb}
\usepackage{enumerate}
\newtheorem{theorem}{Theorem}
\title{Equivariant formality and representation theory}
\author{Chi-Kwong Fok}
\date{Source: arXiv:2609.00448v1 (CC BY 4.0)}
\begin{document}
\maketitle

\noindent\emph{Source and licence.} Equivariant formality and representation theory. Version arXiv:2609.00448v1 (CC BY 4.0), distributed by arXiv under the Creative Commons Attribution 4.0 International licence, http://creativecommons.org/licenses/by/4.0/, as recorded in the arXiv licence metadata for this exact version and shown on the abstract page https://arxiv.org/abs/2609.00448v1. Full licence notice: \texttt{licenses/arxiv-2609.00448-CC-BY-4.0.txt}.\par\smallskip

\noindent\textit{Editorial scope.} Counted unit: the article's theorem circle (source0.tex lines 104--106). The article's complete original proof of it is delivered with it (source0.tex lines 540--552, one \texttt{proof} environment of 13 lines, which the article places away from the statement). No mathematics is written, repaired or re-derived: the statement and the proof are the article's own lines, in the article's own notation, and no retained line is substituted or re-flowed.  No further source-local context is needed: every cross-reference of every retained line resolves inside the two delivered spans.  Nothing was omitted from any retained span, so the package carries no omission record.  No retained line contains a citation, so the wrapper carries no bibliography and no bibliography entry.  No retained line contains a figure environment, an image-inclusion command or drawing code, so no image is delivered and the package declares no asset dependency.  Editorial formatting only, in addition: an AMS \texttt{amsart} preamble that restates the article's own declarations and macros used by the retained lines, namely the environments \texttt{theorem} on separate counters, together with the standard packages those lines need.  The retained environments print in this document as Theorem 1 in order of appearance, whereas the article prints the same items with the numbering of its own full document.  The complete native source of the article as distributed in the arXiv e-print for this exact version is bundled unchanged as \texttt{sources/209028/source0.tex}.
\begin{theorem}\label{circle}
	Let $S$ be a circular subgroup of a compact connected semisimple Lie group $G$. Then $(G, S)$ is isotropy formal if and only if any representation of $G$, when restricted to $S$, is self-dual. 
\end{theorem}

\begin{proof}[Proof of Theorem \ref{circle}] Let $\rho_1, \rho_2, \cdots, \rho_m$ be representations which generate $R(G)$, and $\displaystyle i^*\rho_k=\sum_{n}a_{k, n}t^n\in R(S)$, where $t$ is one-dimensional standard representation of $S$. Suppose that any representation of $G$, when restricted to $S$, is self-dual. Then $a_{k, n}=a_{k, -n}$ for all $k$ and $n$. If we denote $t+t^{-1}-2$ by $s$, then $i^*\rho_k$ can be rewritten as a polynomial $p_k(s)$ of $s$. The regularity of $R:=i^*R(G; \mathbb{Q})$ at the augmentation ideal $i^*I(G; \mathbb{Q})=R\cap (t-1)$ amounts to the smoothness of $\text{Spec} R$ at the point $i^*I(G; \mathbb{Q})$. This is also equivalent to the smoothness of the image of the map $i_\mathbb{R}: \text{Spec} R(S; \mathbb{R})=\text{Spec}\mathbb{R}[t, t^{-1}]\to \text{Spec} R(G; \mathbb{R})=\text{Spec}\mathbb{R}[x_1, x_2, \cdots, x_{m}]/\mathcal{I}\subseteq\text{Spec}\mathbb{R}[x_1, x_2, \cdots, x_m]$ at the point $i_\mathbb{R}((t-1))=I(G; \mathbb{R})$ where $\mathcal{I}$ is the kernel of the map $\mathbb{R}[x_1, x_2, \cdots, x_m]\to R(G; \mathbb{R})$ which sends $x_i$ to $\rho_i$. In other words, the curve with the parametrization $s\mapsto (p_1(s), p_2(s), \cdots, p_m(s))\in\mathbb{R}^m$, which is a reparametrization of the curve $\displaystyle t\mapsto \left(\sum_n a_{1, n}t^n, \cdots, \sum_n a_{m, n}t^n\right)$, is smooth at the point $s=0$ (note that $t=1$ if and only if $s=0$). Thus it suffices to show that $(p_1'(0), p_2'(0), \cdots, p_m'(0))\neq 0$. Differentiating the equation $\displaystyle p_k(s)=\sum_n a_{k, n}t^n$ with respect to $t$ and applying chain rule to the LHS, we get $p_k'(0)$ as follows.
\begin{align*}
	\frac{d}{ds}p_k(s)\cdot\frac{ds}{dt}&=\sum_n na_{k, n}t^{n-1}\\
	\frac{d}{ds}p_k(s)\cdot (1-t^{-2})&=\sum_{n>0}na_{k, n}(t^{n-1}-t^{-n-1})\\
	\frac{d}{ds}p_k(s)&=\sum_{n>0}na_{k, n}(t^{n-1}+t^{n-3}+\cdots+t^{1-n})\\
	p_k'(0)&=\sum_{n>0}n^2a_{k, n}
\end{align*}
The RHS of the last line is not zero because $a_{k, n}\geq 0$ but there exists $n=n_0>0$ such that $a_{k, n_0}>0$. This shows that $R$ is regular at $i^*I(G; \mathbb{Q})$. 

Now suppose that there exists a representation of $G$ whose restriction to $S$ is not self-dual. Let $U$ be a set of representations of $G$ which generate $R(G; \mathbb{R})$, $\overline{U}$ the set of dual representations, and $U\cup \overline{U}=\{\rho_1, \rho_2, \cdots, \rho_m\}$. In this way $\{\rho_1, \rho_2, \cdots, \rho_m\}$ generates $R(G; \mathbb{Q})$ and has the nontrivial involution of taking dual (If $G$ is in addition simply-connected, the set $(\rho_1, \rho_2, \cdots, \rho_m)$ can be taken to be the set of fundamental representations as it generates $R(G; \mathbb{Q})$ and is invariant under the involution of taking dual). Let this involution send $\rho_k$ to $\rho_{\sigma(k)}$, where $\sigma$ is a nontrivial involution on $\{1, 2, \cdots, m\}$. We want to show that $R$ is not regular at $i^*I(G; \mathbb{Q})$. We have that $i^*\rho_k=q_k(t)$ for some Laurent polynomial $q_k$. By the same reasoning as in the previous paragraph, it suffices to show that the image curve $t\mapsto (q_1(t), q_2(t), \cdots, q_m(t))\subseteq \mathbb{R}^m$ for $t>0$ is not smooth at the point $t=1$. Suppose for the sake of contradiction that the curve is smooth at $t=1$. We shall first show that the map $t\mapsto (q_1(t), q_2(t), \cdots, q_m(t))$ is injective for $t$ in an open neighborhood of 1. Choose $k$ such that $k\neq \sigma(k)$. It suffices to show that the map $t\mapsto(q_k(t), q_{\sigma(k)}(t))$ is injective in an open neighborhood of 1. As $q_k(t)$ is a nonconstant Laurent polynomial, there exists $\delta_1>0$ such that $q'_k(t)>0$ or $q'_k(t)<0$ for $t\in(1, 1+\delta_1)$. So $q_k$ and hence $t\mapsto(q_k(t), q_{\sigma(k)}(t))$ is injective for $t\in(1, 1+\delta_1)$. Similarly, there exists $\delta_2>0$ such that $q_k(t)-q_{\sigma(k)}(t)>0$ (resp. $q_k(t)-q_{\sigma(k)}(t)<0$) for $t\in(1, 1+\delta_2)$. Take $\delta=\min\{\delta_1, \delta_2\}$. For $\displaystyle t\in\left(\frac{1}{1+\delta}, 1\right)$, $t^{-1}\in(1, 1+\delta)$, and the map $t\mapsto(q_k(t), q_{\sigma(k)}(t))=(q_{\sigma(k)}(t^{-1}), q_k(t^{-1}))$ is injective and $q_{\sigma(k)}(t^{-1})-q_k(t^{-1})<0$ (resp. $q_{\sigma(k)}(t^{-1})-q_k(t^{-1})>0$). Thus $t\mapsto (q_k(t), q_{\sigma(k)}(t))$ is injective for $\displaystyle t\in\left(\frac{1}{1+\delta}, 1+\delta\right)$. We may further restrict $\displaystyle\left(\frac{1}{1+\delta}, 1+\delta\right)$ to a smaller neighborhood of 1 so that the part of the curve for $t$ in that neighborhood is smooth. Now we can reparametrize that part of the curve by $(r_1(s), r_2(s), \cdots, r_m(s))$ so that the tangent vectors with respect to $s$ is nowhere vanishing (we may parametrize by arc length, for example). In particular the tangent vector at the point $(q_1(1), q_2(1), \cdots, q_m(1))=(\text{dim}\rho_1, \text{dim}\rho_2, \cdots, \text{dim}\rho_m)$ is nonzero. Let $s=s_0$ correspond to the point $(\text{dim}\rho_1, \text{dim}\rho_2, \cdots, \text{dim}\rho_m)$. Note that, since $G$ is semisimple, for any $k$ $q_k(t)$ is of the form $\displaystyle\sum_n a_{k, n}t^n$ where $\displaystyle\sum_n na_{k, n}=0$ and $a_{k, n}\geq 0$. By the AM-GM inequality, 
\[\displaystyle\sum_n a_{k, n}t^n\geq\left(\sum_n a_{k, n}\right)\sqrt{\prod_n (t^n)^{a_{k, n}}}=\text{dim}\rho_k,\] 
and equality holds if and only if $t=1$. Thus each coordinate $r_k(s)$ of the curve achieves its minimum at $s=s_0$ and so $r'(s_0)=(r'_1(s_0), r'_2(s_0), \cdots, r'_m(s_0))=(0, 0, \cdots, 0)$, contradicting the nonvanishing of the tangent vector at $(\text{dim}\rho_1, \text{dim}\rho_2, \cdots, \text{dim}\rho_m)$. It follows that the curve is not smooth at the point $(\text{dim}\rho_1, \text{dim}\rho_2, \cdots, \text{dim}\rho_m)$, which more precisely is a cusp given that each coordinate of the curve is bounded below by the corresponding coordinate of the point.
\end{proof}

\end{document}
